%-----------------------------------------------------------------------------------------------------------------------------------------------%
%	The MIT License (MIT)
%
%	Copyright (c) 2015 Jan Küster
%
%	Permission is hereby granted, free of charge, to any person obtaining a copy
%	of this software and associated documentation files (the "Software"), to deal
%	in the Software without restriction, including without limitation the rights
%	to use, copy, modify, merge, publish, distribute, sublicense, and/or sell
%	copies of the Software, and to permit persons to whom the Software is
%	furnished to do so, subject to the following conditions:
%	
%	THE SOFTWARE IS PROVIDED "AS IS", WITHOUT WARRANTY OF ANY KIND, EXPRESS OR
%	IMPLIED, INCLUDING BUT NOT LIMITED TO THE WARRANTIES OF MERCHANTABILITY,
%	FITNESS FOR A PARTICULAR PURPOSE AND NONINFRINGEMENT. IN NO EVENT SHALL THE
%	AUTHORS OR COPYRIGHT HOLDERS BE LIABLE FOR ANY CLAIM, DAMAGES OR OTHER
%	LIABILITY, WHETHER IN AN ACTION OF CONTRACT, TORT OR OTHERWISE, ARISING FROM,
%	OUT OF OR IN CONNECTION WITH THE SOFTWARE OR THE USE OR OTHER DEALINGS IN
%	THE SOFTWARE.
%	
%
%-----------------------------------------------------------------------------------------------------------------------------------------------%


%============================================================================%
%
%	DOCUMENT DEFINITION
%
%============================================================================%

%we use article class because we want to fully customize the page and dont use a cv template
\documentclass[10pt,A4]{article}	


%----------------------------------------------------------------------------------------
%	ENCODING
%----------------------------------------------------------------------------------------

%we use utf8 since we want to build from any machine
\usepackage[utf8]{inputenc}		

%----------------------------------------------------------------------------------------
%	LOGIC
%----------------------------------------------------------------------------------------

% provides \isempty test
\usepackage{xifthen}

%----------------------------------------------------------------------------------------
%	FONT
%----------------------------------------------------------------------------------------

% some tex-live fonts - choose your own

%\usepackage[defaultsans]{droidsans}
%\usepackage[default]{comfortaa}
%\usepackage{cmbright}
\usepackage[default]{raleway}
%\usepackage{fetamont}
%\usepackage[default]{gillius}
%\usepackage[light,math]{iwona}
%\usepackage[thin]{roboto} 

% set font default
\renewcommand*\familydefault{\sfdefault} 	
\usepackage[T1]{fontenc}

% more font size definitions
\usepackage{moresize}		


%----------------------------------------------------------------------------------------
%	PAGE LAYOUT  DEFINITIONS
%----------------------------------------------------------------------------------------

%debug page outer frames
%\usepackage{showframe}			


%define page styles using geometry
\usepackage[a4paper]{geometry}		

% for example, change the margins to 2 inches all round
\geometry{top=1.25cm, bottom=-.6cm, left=1.5cm, right=1.5cm} 	

%use customized header
\usepackage{fancyhdr}				
\pagestyle{fancy}

%less space between header and content
\setlength{\headheight}{-5pt}		


%customize entries left, center and right
%\lhead{}
%\chead{}
%\rhead{}


%indentation is zero
\setlength{\parindent}{0mm}

%----------------------------------------------------------------------------------------
%	TABLE /ARRAY DEFINITIONS
%---------------------------------------------------------------------------------------- 

%for layouting tables
\usepackage{multicol}			
\usepackage{multirow}

%extended aligning of tabular cells
\usepackage{array}

\newcolumntype{x}[1]{%
>{\raggedleft\hspace{0pt}}p{#1}}%


%----------------------------------------------------------------------------------------
%	GRAPHICS DEFINITIONS
%---------------------------------------------------------------------------------------- 

%for header image
\usepackage{graphicx}

%for floating figures
\usepackage{wrapfig}
\usepackage{float}
%\floatstyle{boxed} 
%\restylefloat{figure}

%for drawing graphics		
\usepackage{tikz}				
\usetikzlibrary{shapes, backgrounds,mindmap, trees}


%----------------------------------------------------------------------------------------
%	Color DEFINITIONS
%---------------------------------------------------------------------------------------- 

\usepackage{color}

%accent color
\definecolor{sectcol}{RGB}{0,150,255}

%dark background color
\definecolor{bgcol}{RGB}{110,110,110}

%light background / accent color
\definecolor{softcol}{RGB}{225,225,225}


%============================================================================%
%
%
%	DEFINITIONS
%
%
%============================================================================%

%----------------------------------------------------------------------------------------
% 	HEADER
%----------------------------------------------------------------------------------------

% remove top header line
\renewcommand{\headrulewidth}{0pt} 

%remove botttom header line
\renewcommand{\footrulewidth}{0pt}	  	

%remove pagenum
\renewcommand{\thepage}{}	

%remove section num		
\renewcommand{\thesection}{}			

%----------------------------------------------------------------------------------------
% 	ARROW GRAPHICS in Tikz
%----------------------------------------------------------------------------------------

% a six pointed arrow poiting to the left
\newcommand{\tzlarrow}{(0,0) -- (0.2,0) -- (0.3,0.2) -- (0.2,0.4) -- (0,0.4) -- (0.1,0.2) -- cycle;}	

% include the left arrow into a tikz picture
% param1: fill color
%
\newcommand{\larrow}[1]
{\begin{tikzpicture}[scale=0.58]
	 \filldraw[fill=#1!100,draw=#1!100!black]  \tzlarrow
 \end{tikzpicture}
}

% a six pointed arrow poiting to the right
\newcommand{\tzrarrow}{ (0,0.2) -- (0.1,0) -- (0.3,0) -- (0.2,0.2) -- (0.3,0.4) -- (0.1,0.4) -- cycle;}

% include the right arrow into a tikz picture
% param1: fill color
%
\newcommand{\rarrow}
{\begin{tikzpicture}[scale=0.7]
	\filldraw[fill=sectcol!100,draw=sectcol!100!black] \tzrarrow
 \end{tikzpicture}
}



%----------------------------------------------------------------------------------------
%	custom sections
%----------------------------------------------------------------------------------------

% create a coloured box with arrow and title as cv section headline
% param 1: section title
%
\newcommand{\cvsection}[1]
{
	\begin{center}
		\large\textcolor{sectcol}{\textbf{#1}}
	\end{center}
}

%create a coloured arrow with title as cv meta section section
% param 1: meta section title
%
\newcommand{\metasection}[2]
{
%\begin{tabular*}{1\textwidth}{r r}
\footnotesize{#2} \hspace*{\fill} \footnotesize{#1}\\[1pt]
%\end{tabular*}
}

%----------------------------------------------------------------------------------------
%	 CV EVENT
%----------------------------------------------------------------------------------------

% creates a stretched box as cv entry headline followed by some paragraphs about 
% the work you did
% param 1:	event time i.e. 2014 or 2011-2014 etc.
% param 2:	event name (what did you do?)
% param 3:	institution (where did you work / study)
% param 4:	list of paragraphs
%
\newcommand{\cvevent}[4]
{

\begin{tabular*}{1\textwidth}{p{13.6cm}  x{3.9cm}}
	\textbf{#2} - \textcolor{bgcol}{#3} &   \vspace{2.5pt}\textcolor{sectcol}{#1}
\end{tabular*}

\vspace{-8pt}
\textcolor{softcol}{\hrule}
\vspace{6pt}

	\foreach \desc in {#4}{
		$\cdot$ \desc\\[3pt]
	}
	
\vspace{3pt}
}

% creates a stretched box as 
\newcommand{\cveventmeta}[2]
{
	\mbox{\mystrut \hspace{87pt}\textit{#1}}\\
	#2
}

%----------------------------------------------------------------------------------------
% CUSTOM STRUT FOR EMPTY BOXES
%----------------------------------------- -----------------------------------------------
\newcommand{\mystrut}{\rule[-.3\baselineskip]{0pt}{\baselineskip}}

%----------------------------------------------------------------------------------------
% CUSTOM LOREM IPSUM
%----------------------------------------------------------------------------------------
\newcommand{\lorem}
{Lorem ipsum dolor sit amet, consectetur adipiscing elit. Donec a diam lectus.}



%============================================================================%
%
%
%
%	DOCUMENT CONTENT
%
%
%
%============================================================================%
\begin{document}


%use our custom fancy header definitions
\pagestyle{fancy}	





%----------------------------------------------------------------------------------------
%	HEADER IMAGE
%----------------------------------------------------------------------------------------

%\begin{figure}[H]
%\begin{flushright}
%	\includegraphics[trim= 320 130 460 210,clip,width=0.2\linewidth]{myfoto.jpg}	%trimming relative to image size!
%\end{flushright}
%\end{figure}


%---------------------------------------------------------------------------------------
%	TITLE HEADLINE
%----------------------------------------------------------------------------------------
\vspace{-8pt}
\begin{center}
	\Huge \textsc{Antonin Cholet} \textcolor{sectcol}{\rule[-1mm]{1mm}{0.9cm}} \textsc{C.V.}\\[2pt]
	\small Avocat / DPO / Webdev
\end{center}



\vspace{6pt}


%---------------------------------------------------------------------------------------
%	META SECTION
%----------------------------------------------------------------------------------------
\metasection{www.cholet-avocat.fr}{\textbf{Status:} Avocat libéral}
\metasection{adresse courriel}{\textbf{Langues} Anglais (Niveau C1 TOEIC 825/990, février 2010), Italien (Niveau C1)}
\metasection{+33 (0)767250039}{\textbf{Techs:} Ruby, Javascript, Bash, Git, Pandoc, Latex, HTML/CSS, Markdown}
\vspace{-2pt}
\textcolor{softcol}{\hrule}
\vspace{6pt}

\normalsize

%---------------------------------------------------------------------------------------
%	SUMMARAY (optional)
%----------------------------------------------------------------------------------------
\vspace{-6pt}
\cvsection{Préambule}
Avocat en droit public et droit des données personnelles avec des compétences en développement web Fullstack Ruby, ancien DAF-DRH d’un établissement public local, et ancien responsable contentieux d’une URSSAF, ancien juriste dans une entreprise publique.\\


%============================================================================%
%
%	CV SECTIONS AND EVENTS (MAIN CONTENT)
%
%============================================================================%

%---------------------------------------------------------------------------------------
%	EXPERIENCE
%----------------------------------------------------------------------------------------
\cvsection{Experience}

\cvevent{2017 - ce jour}{Avocat en droit public}{En libéral (Besançon)}{
  {Conseil et contentieux en Urbanisme / Fonction publique / Collectivités locales / Données personnelles},{Représentation devant les juridictions administratives}}

\cvevent{2017}{Coordonnateur réglementation et sécurisation juridique}{URSSAF (Besançon)}{
  {Management du service contentieux / Gestion des contentieux sociaux en matière de recouvrement}}

\cvevent{2015 - 2016}{Daf - drh}{Sivuresc (Blanc Mesnil)}{
  {Budget en M14 (plus de 7 millions d'euros) / Mandatements et titres / Comptes administratifs et de gestion de fin d’exercice / Coordination avec le comptable public},
  {Conseil syndical / Marchés publics / Controle de légalité \& rescrits fiscaux},
  {Planification RH / Réformes (RIFSEEP, contractuels) / Négociations syndicales / Évaluations \& formations / Carrières \& payes},
  {Management et réorganisation d’un service administratif}
}

\cvevent{2013 - 2015}{Avocat en droit public}{Lazare Avocats (Paris)}{
	{En collaboration : Urbanisme / Préemption \& expropriation / Collectivités locales / Propriétés publiques},
	{En libéral : Fonction publique},
  {Représentation devant les juridictions administratives}
}

\cvevent{2012 - 2013}{Avocat en stage}{Haute école des avocats conseils (Versailles)}{
	{Selarl Soler-Couteaux/Llorens (Strasbourg) - \textit{auprès de Me Soler-Couteaux, avocat et professeur agrégé} : Urbanisme / Propriétés publiques},
	{Tribunal administratif (Rouen) - \textit{auprès de M. AUPOIX, président 1ère chambre} : Notes de rapporteur / projets de jugement et d'ordonnance},
  {A5 Avocats Associés (Paris) - \textit{auprès de Me Cayla-Destrem} : Fonction publique / Collectivités locales}
}

\cvevent{2010 - 2011}{Juriste}{Enedis (La Defense)}{
  {Appui juridique opérationnel aux directions régionales et agents en droit de l'Énergie / Civil / Urbanisme / Environnement}
}

\cvevent{2010}{Juriste en stage}{Université de Franche Comté (Besançon)}{
  {Adamas Avocats (Lyon) - \textit{auprès de Me SANTY} : Conseil et rédaction d'acte en droit des sociétés / contrats},
  {Fidal (Besançon)- \textit{auprès de Me HOURNON} : Conseil en droit de la concurrence / droit des contrats}
}




%---------------------------------------------------------------------------------------
%	EDUCATION SECTION
%--------------------------------------------------------------------------------------
\cvsection{Formation initiale}

\cvevent{2013}{Capacité d'aptitude à la profession d'avocat (CAPA)}{HEDAC de Versailles}{}

%\textcolor{softcol}{\hrule}

%
\cvevent{2011}{Master II Droit de l’Environnement}{Université de Versailles - Saint Quentin en Yvelines}{}

%\textcolor{softcol}{\hrule}

%
\cvevent{2010}{Master II Juriste d’affaires européen}{Université de Franche Comté}{}

%\textcolor{softcol}{\hrule}

%
\cvevent{2007 - 2008}{Erasmus}{Università di Roma Tre}{}

\cvevent{2004}{Bac S}{Lycée Pasteur (Besançon)}{}

\cvsection{Formation continue}

\cvevent{2021}{Learn command line}{Codecademy}{}

\cvevent{2021}{Learn bash scripting}{Codecademy}{}

\cvevent{2019 - 2021}{Certificat de spécialisation de Délégué à la protection des données}{CNAM (Paris)}{
  {Coordinateur de formation : M. Desgens-Pasanau, magistrat, professeur des universités},
  {Droit de la protection des données / Droit du numérique / Métier du DPO}
}

\cvevent{2020}{Certificat de Concepteur - développeur d’applications web Fullstack}{LE WAGON}{
  {Ruby \& Rails / MVC Pattern / JavaScript},
  {Database \& SQL / API calls},
  {HTML \& CSS / SCSS / Bootstrap},
  {Git / Terminal},
  {Projet final d’une app web réutilisant les données publiques  et géodonnées publiques d’urbanisme, du cadastre, et des adresses (par requetes API sur apicarto.ign.fr et adresse.data.gouv.fr)}
}

\cvevent{2020}{Learn ruby}{Codecademy}{}

\cvevent{2019}{MOOC Protection des données personnelles}{CNAM}{}

\cvevent{2018}{Actualité du droit de l’urbanisme - Me Soler-Couteaux}{ERAGE}{}

\cvevent{2018}{Responsabilité civile professionnelle de l’avocat}{ERAGE}{}

\cvevent{2018}{La médiation}{Lexbase formation}{}

\cvevent{2018}{Réparation du préjudice corporel}{Fides Formation}{}

\cvevent{2016}{Le management des ressources humaines}{CNFPT}{}

\cvevent{2016}{Animation, direction et mobilisation de son service}{CNFPT}{}

\cvevent{2015}{La préparation, l'élaboration et l'exécution du budget M14}{CNFPT}{}

\cvevent{2015}{Les bases des finances locales pour les cadres}{CNFPT}{}

%--------------------------------------------------------------------------------------------------
%	ARTIFICIAL FOOTER (fancy footer cannot exceed linewidth) 
%--------------------------------------------------------------------------------------------------

% \null
% \vspace*{\fill}
% \hspace{-0.25\linewidth}\colorbox{white}{\makebox[1.5\linewidth][c]{\mystrut  \textnormal{\textcolor{sectcol}{www.jankuester.com} $\cdot$ \textcolor{sectcol}{github.com/jankapunkt}}}}




%============================================================================%
%
%
%
%	DOCUMENT END
%
%
%
%============================================================================%
\end{document}
